\documentclass[11pt]{article}
\usepackage{amsmath,amssymb,amsthm}
\usepackage[margin=1.1in]{geometry}
\usepackage{hyperref}

\newtheorem{theorem}{Theorem}
\newtheorem{conjecture}{Conjecture}
\newtheorem{lemma}{Lemma}
\theoremstyle{definition}
\newtheorem{measurement}{Measurement}

\title{An Empirical Map of the Goldbach Conditional Chain:\\
from a Sub-Poissonian Constant to a Direct Portrait\\
of the M\"obius-Twisted Elliott--Halberstam Axiom}
\author{[author]\thanks{AI-assisted exploratory program, 2026.
Code and full logs: \url{https://github.com/soke91/goldbach-ln2}}}
\date{Draft — \today}

\begin{document}
\maketitle

\begin{abstract}
By a theorem of Huang and Li (continuing Pan), the binary Goldbach
conjecture for large even $N$ follows from the Bombieri--Vinogradov
theorem together with a single unproven statement: the M\"obius-twisted
Elliott--Halberstam hypothesis $EH_\mu(N^{\theta'})$ for some
$\theta' > 1/2$, needed only for the fixed residue class $N \bmod q$,
whose object is the correlation sequence $c(n) = \Lambda(n)\mu(N-n)$.
We present a systematic empirical program measuring this landscape and
everything feeding it: (i) the fixed-residue discrepancy of $c(n)$ is
random-walk-sized straight through the $\sqrt N$ barrier
($\theta = 0.30 \to 0.70$, two scales, five values of $N$);
(ii) the innermost sum $T(N) = \sum_n \Lambda(n)\mu(N-n)$ --- whose
$o(N)$ bound is itself beyond unconditional technology --- follows a
textbook half-normal law over hundreds of values of $N$; (iii) sifted
prime counts sit on the Buchstab curve $e^\gamma\omega(u)$ to
$\pm 2\cdot 10^{-4}$ throughout the sieve-blind zone (``Conjecture P'');
(iv) an exact structure law
$s(n) = s_0(x)\prod_{q \mid n,\, q > y}(q-1)/(q-2)$ holds to 3--4
decimals at roughness $x^{1/8}$ across three scales; and (v) the
fluctuation of Goldbach counts against the Hardy--Littlewood prediction
carries a canonical dispersion constant equal to $\ln 2$ within
$0.9\%$, an invariant of additive prime correlations of order $\ge 3$.
Twenty-four documented self-corrections and three pre-registered blind
extrapolation tests (all hits) accompany the program. No theorem toward
Goldbach is claimed; the contribution is a reproducible, falsifiable
measurement corpus locating the entire parity obstruction in one named
technical question.
\end{abstract}

\section{Introduction}
% Frame: Chen -> parity -> Huang-Li reduction; what measurement can and
% cannot contribute; honesty protocol (corrections, blind tests).
[Draft: import from RESEARCH\_NOTE.md intro + Closure paragraphs.]

\section{The conditional chain and the final axiom}
% Huang-Li Theorem 1 / Corollary 1; Pan decomposition f_0..f_3;
% how EH_mu is consumed (mu(d) log d weighted modulus sums).
[Draft: import from FINAL\_BOSS \S5--\S7 (translated).]

\section{Direct measurements of the axiom's landscape}
\subsection{Fixed-residue discrepancy of $c(n)$ through the barrier}
% Tables: theta 0.30-0.70, N at 1e8 (x3) and 1e9 (x2); R_c stats.
\subsection{The innermost sum $T(N)$}
% 220 N at 1e8 (+ ladder at 1e9): half-normal statistics.
\subsection{Smooth (well-factorable-type) moduli}
% 479 squarefree smooth moduli: equal health.
\subsection{Plain M\"obius landscape to $\theta = 0.88$}

\section{Conjecture P: sifted primes on the Buchstab curve}
% P-profile: R_P = 1.0000(2) in the blind zone; fine structure =
% omega oscillation; 1/log z finite-size law; scale stability.

\section{Structure laws at the target roughness}
% theta = 1/8 exact law; worst class; chi^2 fiber uniformity ladder
% (1e6 -> 1e12); pre-registered test killing Poisson convergence at 26 sigma.

\section{The $\ln 2$ dispersion constant}
% Canonical estimator; 18-octave fitted limit 0.6931(61); three
% independent constructions; gap-permutation kill test (>= 3-point).

\section{Loss anatomy and the minimal target}
% U_req(theta) profile; loss integral 4.08 vs 1.26; P_loc: two
% aggregates per n, factor-125 slack; positivity-propagation remarks.

\section{Methods, corrections, and falsification}
% 24 corrections list; blind-test protocol; window systematics saga;
% reproducibility (script table).

\section*{Acknowledgements}
% Huang-Li, Pan, and the literature that anchors the frame.

\begin{thebibliography}{9}
\bibitem{HL} J.-J.~Huang, H.~Li, \emph{On the connection between the
Goldbach conjecture and the Elliott--Halberstam conjecture},
arXiv:2005.03811; in: Combinatorial and Additive Number Theory IV,
Springer (2021).
\bibitem{Pan} C.~D.~Pan, \emph{A new attempt on Goldbach conjecture},
Chin.\ Ann.\ Math.\ 3 (1982).
\bibitem{MV} M.~R.~Murty, A.~Vatwani, \emph{Twin primes and the parity
problem}, J.\ Number Theory 180 (2017).
\bibitem{Chen} J.~R.~Chen, \emph{On the representation of a larger even
integer as the sum of a prime and the product of at most two primes},
Sci.\ Sinica 16 (1973).
\end{thebibliography}

\end{document}
